\section{F-SM-31: the full interval hypotheses do not force AB5 or presentability}

\subsection{Statement}
For every field \(k\), there is a complete and cocomplete
\(k\)-linear abelian symmetric monoidal category \(\mathcal M\)
with enough projectives satisfying all the precise adjoint,
unit-normalization, and full tensor-interval conditions of
cote 134-8, footers 47--48, which is neither AB5 nor locally
presentable.  Its tensor product is exact and preserves colimits
in both variables,
and the scalar functor is exact and strong symmetric monoidal.
Thus the missing Grothendieck hypothesis in F-SM-30 cannot be
deduced from those original conditions.

This is a counterexample to that proposed implication, not to
the original localization conclusion.  The latter holds in the
example, as is proved at the end.  No double-dual, ultrapower,
or unreconstructed tensor-coherence assertion is used.

\subsection{Construction and complete coherence check}
Let \(\mathcal V=\operatorname{Vect}_k\) and set
\[
 \mathcal M=\operatorname{Ch}_{\ge0}(k)\times\mathcal V^{\rm op}.
                                                               \tag{SM31.1}
\]
We denote its objects by \((C,A)\), where \(A\) is an ordinary
vector space but belongs to the opposite-category factor.
Thus a morphism \((C,A)\to(D,B)\) consists of a chain map
\(f:C\to D\) and a linear map \(g:B\to A\).
The category is abelian, complete and cocomplete, componentwise.
Every object of \(\mathcal V^{\rm op}\) is projective because
every monomorphism of vector spaces splits.  The chain category
has the usual projective generators \(k[0],D^n(k)\), \(n\ge1\).
Consequently \(\mathcal M\) has enough projectives.
It even has a projective generator: combine those chain
generators with \((0,k)\), and take their coproduct.
The latter generates the opposite factor because linear
functionals to \(k\) separate nonzero maps of vector spaces.

For a vector space \(U\), the canonical scalar action on
\(\mathcal V^{\rm op}\) is
\[
             U\odot A=\operatorname{Hom}_k(U,A),
             \quad\text{regarded as an object of }\mathcal V^{\rm op}.
                                                               \tag{SM31.2}
\]
It is covariant in \(U\) after the opposite-category convention.
The canonical currying isomorphism gives
\((U\otimes V)\odot A\simeq U\odot(V\odot A)\),
and \(k\odot A\simeq A\).
Define the tensor product by
\[
 (C,A)\boxtimes(D,B)=
 \left(C\otimes_kD,
    \operatorname{Hom}_k(C_0,B)\oplus
    \operatorname{Hom}_k(D_0,A)\right).                    \tag{SM31.3}
\]
The finite sum in the second component is the same vector
space for products and coproducts.  For morphisms
\((f,g):(C,A)\to(C',A')\) and
\((h,j):(D,B)\to(D',B')\), the map on this component,
in the ordinary vector-space direction, is
\[
 \begin{aligned}
 \operatorname{Hom}(C'_0,B')&\longrightarrow\operatorname{Hom}(C_0,B),
                 &q&\longmapsto j\,q\,f_0,\\
 \operatorname{Hom}(D'_0,A')&\longrightarrow\operatorname{Hom}(D_0,A),
                 &q&\longmapsto g\,q\,h_0.
 \end{aligned}                                             \tag{SM31.4}
\]
This defines a bifunctor additive in each variable.
Its unit is \(\mathbf1=(k[0],0)\).
In fact the tensor is exact in each variable: evaluation at
chain degree zero is exact, all vector-space short exact
sequences split, and tensor and either variance of Hom over
a field preserve their exactness, with the opposite variance
accounted for as in (SM31.4).  It preserves coproducts as well.
For example, their second component satisfies
\[
 \operatorname{Hom}\Bigl(\bigoplus_i(C_i)_0,B\Bigr)
 \oplus\operatorname{Hom}\Bigl(D_0,\prod_iA_i\Bigr)
 \simeq\prod_i\bigl(\operatorname{Hom}((C_i)_0,B)
                 \oplus\operatorname{Hom}(D_0,A_i)\bigr).
\]
This is the required coproduct identity in the opposite
factor.  Since any colimit in an abelian category is obtained
from coproducts and a cokernel, tensor preserves all colimits
in each variable.  In particular the example also satisfies
the stronger right-exactness conventions for a tensor product.

We specify and check all coherence maps.  In the first component
use the standard associator, unit maps, and signed symmetry of
nonnegative chain complexes.  In the second component an
\(m\)-fold tensor of \((C_i,A_i)\), in any parenthesization,
has the canonical normal form
\[
 \bigoplus_{i=1}^m
  \operatorname{Hom}_k\left(
      \bigotimes_{j\ne i}(C_j)_0,\ A_i\right),             \tag{SM31.5}
\]
where the tensor factors are in their original order.
This follows by induction from (SM31.3), distributivity over
finite sums, and currying.  Define every associator by the
canonical identifications with this normal form.  The two
pentagon composites coincide because on each labelled summand
they are the same currying map on the same ordered tensor
factors.  Define the symmetry by swapping the two components
in (SM31.3) and using the usual interchange of degree-zero
tensor factors in (SM31.5).  The hexagon composites induce the
same permutation of these factors and the same permutation
of the summand labels, and hence coincide.  The triangle and
unit axioms reduce to \(\operatorname{Hom}_k(k,A)=A\).
The first-component signed symmetry contributes no sign to
the second component, which uses only chain degree zero.
These identifications are natural, by (SM31.4).
Thus this is a symmetric monoidal category with all the
required coherence, not merely a formula on objects.

\subsection{Scalar adjoints and derived normalization}
The actual scalar copower of the tensor unit is
\[
 \kappa(U)=(U[0],0).
\]
Indeed the second component of any coproduct of copies of the
zero object remains zero, and scalar presentations give the
displayed first component.  Formula (SM31.3) makes \(\kappa\)
strong symmetric monoidal.  It is exact and preserves all
small limits and colimits, by the componentwise constructions.
Its two adjoints are
\[
 \varphi(C,A)=H_0(C),\qquad
 \psi(C,A)=C_0,
 \qquad \varphi\dashv\kappa\dashv\psi.                    \tag{SM31.6}
\]
For the left adjunction, a chain map \(C\to U[0]\) is exactly
a linear map \(H_0(C)\to U\); for the right adjunction, a
chain map \(U[0]\to C\) is exactly a linear map \(U\to C_0\).
The opposite-vector component contributes only its unique
zero map in either calculation.  In particular
\(\psi=\operatorname{Hom}_{\mathcal M}(\mathbf1,-)\),
as required in the original setup.

Let \(\Lambda(C,A)=C\).  This is an exact functor to
\(\operatorname{Ch}_{\ge0}(k)\), with \(H_0\Lambda=\varphi\).
A projective object of \(\mathcal M\) has projective first
component in the chain category, whose positive homology
vanishes.  The fully proved finite-diagonal lemma BR-13401
therefore gives, naturally,
\[
                     L\varphi(C,A)\simeq C.                \tag{SM31.7}
\]
The canonical map on the unit is \(k[0]\to k[0]\), the
identity.  This verifies the original unit normalization.

\subsection{The complete tensor interval, including degree zero}
Take \(J=(D^1(k),0)\), where \(D^1(k)\) has identity
differential from degree one to degree zero, and let
\(\alpha:\mathbf1\to J\) be the degree-zero inclusion.
For every \(X=(C,A)\), there is the natural identification
\[
                 J\boxtimes X=(D^1(k)\otimes C,A).          \tag{SM31.8}
\]
Under it, \(\alpha\boxtimes1_X\) has second component the
identity of \(A\).  Its first component is the degreewise
monomorphism \(C\to D^1(k)\otimes C\).
In degree \(n\), write the latter tensor complex as
\(C_n\oplus C_{n-1}\).  Its differential and contraction are
\[
 d(x,y)=(d_Cx+y,-d_Cy),\qquad h(x,y)=(0,x).
\]
Direct substitution gives \(dh+hd=1\).
Consequently \(\alpha\boxtimes1_X\) is monic for every
\(X\), and (SM31.7) gives
\[
                    L\varphi(J\boxtimes X)=0
\]
in every degree, including degree zero.  This verifies the
full original interval condition, without a finiteness
restriction on \(C\) or \(A\).

\subsection{Failure of AB5 and of local presentability}
We give explicit arguments for both failures.
Let \(I=\mathbb N\), let \(V=k^{(I)}\), and for a finite
subset \(S\subset I\) put
\(V_S=k^{(I\setminus S)}\subset V\).
There is an inverse system of short exact vector-space
sequences
\[
          0\longrightarrow V_S\longrightarrow V
                  \longrightarrow V/V_S\longrightarrow0.
\]
Its limits have
\(\lim_S V_S=0\), \(\lim_S V=V\), and
\(\lim_S V/V_S=k^I\).
The resulting map \(V\to k^I\) is not surjective, since
the constant family \((1)_{i\in I}\) is not finitely supported.
Passing to opposite categories turns this into a filtered
system of short exact sequences whose colimit is not exact.
Thus \(\mathcal V^{\rm op}\), and hence \(\mathcal M\),
fails AB5.

For local presentability fix any infinite regular cardinal
\(\lambda\) and any nonzero vector space \(A\).
Let \(I\) be a set of cardinality \(\lambda\), and now
index the same tail subspaces \(V_S=k^{(I\setminus S)}\)
by all subsets \(S\subset I\) of cardinality less than
\(\lambda\).  This indexing poset is \(\lambda\)-filtered,
and its colimit in \(\mathcal V^{\rm op}\) is again zero,
because the intersection of the tail subspaces is zero.
Choose \(0\ne a\in A\).  The map \(V\to A\) taking every
basis vector to \(a\) restricts nontrivially to every \(V_S\).
It therefore defines a nonzero element of
\[
 \operatorname*{colim}_{|S|<\lambda}
          \operatorname{Hom}_k(V_S,A),
\]
whereas \(\operatorname{Hom}_{\mathcal V^{\rm op}}(A,0)=0\).
Thus no nonzero object of \(\mathcal V^{\rm op}\) is
\(\lambda\)-presentable, for any infinite regular \(\lambda\).

If \((C,A)\) is \(\lambda\)-presentable in \(\mathcal M\),
apply the preceding filtered diagram in its second factor,
with zero first factor.  The same calculation forces \(A=0\).
Filtered colimits of objects with zero second factor still
have zero second factor.  Hence no family of
\(\lambda\)-presentable objects can generate \(\mathcal M\)
under \(\lambda\)-filtered colimits: for example \((0,k)\)
cannot be so obtained.  Since this holds for every regular
\(\lambda\), \(\mathcal M\) is not locally presentable.

\subsection{The localization conclusion nevertheless holds here}
By (SM31.7), \(W\) consists exactly of the maps whose first
components are quasi-isomorphisms.  Projection
\(\Lambda(C,A)=C\) has a section \(G(C)=(C,0)\), with
\(\Lambda G=\mathrm{id}\).  The natural map
\((C,A)\to(C,0)\) belongs to \(W\), since its first component
is the identity.  These functors and this natural weak-equivalence
map give mutually inverse functors after localization.  Therefore
\[
                W^{-1}\mathcal M\simeq D_{\ge0}(k).
\]
The example consequently rules out only the proposed deduction
of AB5 or presentability from the original predicates.  It does
not decide the bare abstract localization assertion in general.
\(\square\)

\paragraph{Primary-source scope.}
The original predicates checked here are those of
Grothendieck, A. (n.d.), \emph{Pursuing stacks: Linearization
of homotopy types}, cote 134-8, manuscript footers 47--48,
\url{https://grothendieck.umontpellier.fr/134-8.pdf}.
The construction and all checks above are supplied in full;
the derived-functor identification invokes only the already
proved BR-13401 and BR-13402.
